% !TEX program = lualatex
% Auto-generated from YAML (v3) — 14: データサイエンス・AI（7）：AIモデルの評価

\documentclass[aspectratio=43,professionalfonts,handout,10pt]{beamer}
\usepackage{beamer_template}

\newcommand{\LectureNum}{14}
\newcommand{\LectureTitle}{データサイエンス・AI（7）：AIモデルの評価}
\newcommand{\CourseName}{情報リテラシー}
\newcommand{\TextPages}{-}
\newcommand{\NextTopic}{後期復習}
\newcommand{\NextPages}{後期全範囲}

\begin{document}
% --- Slide 1: title ---

\begin{frame}{\CourseName\quad 第\LectureNum 回「\LectureTitle 」}
  {\small \TermName\quad \TeacherName\quad ／\quad 教科書 \TextPages}

  \vfill

  \begin{block}{今日の流れ}
    \begin{enumerate}
      \item 推論とは
      \item 評価指標①：正解率（Accuracy）
      \item 評価指標②：適合率（Precision）
      \item 評価指標③：再現率（Recall）
      \item 評価指標④：混同行列（Confusion Matrix）
      \item 過学習と未学習
    \end{enumerate}
  \end{block}
  \vfill

  \begin{block}{今日の目標}
    \begin{itemize}
      \item AIの学習と推論，評価を行うことができる。
    \end{itemize}
  \end{block}
  \vfill

  \begin{exampleblock}{キーワード}
    \small \textbf{推論（Inference）} ／ \textbf{第12回（MNIST）での推論} ／ \textbf{評価の必要性}
  \end{exampleblock}
\end{frame}

% --- Slide 3: table ---

\begin{frame}{4つの評価指標}
  \centering
  \begin{tabular}{llll}
    \toprule
    指標 & 定義 & 計算式 & 何を見るか \\
    \midrule
    正解率（Accuracy） & 全体で何\%正解したか & 正解数 ÷ 全体の数 & 全体の性能 \\
    適合率（Precision） & 「この数字だ」と予測したとき，実際に正しい割合 & 正しく予測した数 ÷ その数字だと予測した全体 & 予測の信頼性（間違った予測の少なさ） \\
    再現率（Recall） & 実際のその数字を，どれだけ正しく見つけられたか & 正しく予測した数 ÷ 実際のその数字の数 & 見逃しの少なさ \\
    混同行列（Confusion Matrix） & どの数字をどの数字と間違えたか & 行：実際の数字，列：予測した数字 & 間違いのパターン \\
    \bottomrule
  \end{tabular}
\end{frame}

% --- Slide 4: twocol ---

\begin{frame}{適合率と再現率のトレードオフ}
  \begin{columns}[T]
    \begin{column}{0.48\textwidth}
      \begin{block}{適合率（Precision）重視}
        \small
        \begin{itemize}
          \item 予測した側から見る
          \item 「病気だ」と診断したら確実に病気
          \item 誤診を減らしたい場面
          \item 適合率を上げると再現率が下がる
        \end{itemize}
      \end{block}
    \end{column}
    \begin{column}{0.48\textwidth}
      \begin{exampleblock}{再現率（Recall）重視}
        \small
        \begin{itemize}
          \item 実際の側から見る
          \item 病気の人を見逃さない
          \item 見落としが致命的な場面（がん検診など）
          \item 再現率を上げると適合率が下がる
        \end{itemize}
      \end{exampleblock}
    \end{column}
  \end{columns}
\end{frame}

% --- Slide 5: twocol ---

\begin{frame}{過学習と未学習}
  \begin{columns}[T]
    \begin{column}{0.48\textwidth}
      \begin{block}{過学習（Overfitting）}
        \small
        \begin{itemize}
          \item 訓練データに合わせすぎた状態
          \item 訓練データ99\%正解，テストデータ70\%正解
          \item 新しいデータで性能が落ちる
          \item 原因：モデルが複雑すぎる・データが少ない
        \end{itemize}
      \end{block}
    \end{column}
    \begin{column}{0.48\textwidth}
      \begin{exampleblock}{未学習（Underfitting）}
        \small
        \begin{itemize}
          \item 学習が不足した状態
          \item 訓練データ70\%正解，テストデータ68\%正解
          \item そもそも性能が低い
          \item 原因：モデルが単純すぎる・学習が不足
        \end{itemize}
      \end{exampleblock}
    \end{column}
  \end{columns}
\end{frame}

% --- Slide 6: practice ---

\begin{frame}{Colab実習：MNISTモデルの評価}
  {\small 実際にコンピュータで操作してください。}

  \vfill

  \begin{block}{手順}
    \begin{enumerate}
      \item Part 1：第12回の復習 → MNISTモデルの学習を実行
      \item Part 2：評価指標の計算 → 正解率・適合率・再現率・混同行列を確認
      \item 「7」を「1」と間違えやすい，「7」を「8」と間違えやすい，などのパターンを発見する
    \end{enumerate}
  \end{block}
\end{frame}

% --- Slide 7: exercise ---

\begin{frame}{演習（後期後半まとめ）}
  {\small 各自で取り組んでください。}

  \vfill

  \begin{block}{演習1}
    正解率・適合率・再現率の違いを説明せよ
  \end{block}
  \vfill

  \begin{block}{演習2}
    「がん検診AI」を開発するとき，適合率と再現率のどちらを重視すべきか。理由も述べよ
  \end{block}
  \vfill

  \begin{block}{演習3}
    過学習とはどのような状態か。なぜ問題なのか説明せよ
  \end{block}
\end{frame}

% --- チェックリスト ---

\begin{frame}{まとめ・チェックリスト}
  \begin{block}{自己チェック（できたら $\checkmark$ をつけよう）}
    \begin{itemize}
      \item[$\square$] 正解率・適合率・再現率・F値の評価指標を理解した
      \item[$\square$] 混同行列（Confusion Matrix）の見方と使い方を学んだ
      \item[$\square$] 過学習・未学習の問題とその対策を把握した
    \end{itemize}
  \end{block}

  \vfill

  {\small 質問があれば授業後またはTeamsで受け付けます。}
\end{frame}

\end{document}
